\section{Related work}
\label{sec:related}

\subsection{Construction robotics and the role of building information}

Construction robotics has long been framed as an information problem as much
as a machine problem. Early surveys anticipated that sensing and product
models would be needed to compensate for site variability
\cite{vaha2013automation,bock2015future,ardiny2015review}. Later reviews
identified persistent barriers in localization, work preparation, safety,
interoperability, and adaptation to unstructured conditions
\cite{davila2019challenges,melenbrink2020onsite,xiao2022advancements}. The
diversity of the field is relevant to scene modelling: collective fabrication,
on-site mobile systems, prefabrication, and human--robot collaboration impose
different requirements on the digital representation used at runtime
\cite{petersen2019collective,xu2026review,afaq2026steel}. Consequently,
``robot-ready'' cannot denote a single geometric level of detail or one fixed
set of asset attributes.

Reviews focused specifically on BIM--robot interaction show a progression from
one-way export towards lifecycle information flow and closed-loop exchange
\cite{zhang2022bimfar,anane2023interoperability,yue2025interactions,odugu2025review}.
Across this literature, BIM supplies geometry, object identity, material and
spatial semantics, schedules, and task constraints. The receiving robotic
system commonly creates its own reduced representation for a particular
purpose. A navigation map, an assembly graph, an SDF world, and an articulated
physics asset are all robot-relevant, but they make different claims about
what information is complete and valid. This fragmentation motivates the
present emphasis on an explicit scene contract rather than another destination
format.

For robotic assembly, Dawod and Hanna used BIM geometry and design context to
support object recognition, grasp search, and motion planning for discrete
structures \cite{dawod2019recognition}. Ding et al.\ connected image-based
reconstruction to BIM task-level planning for brick assembly
\cite{ding2020brick}. Wong Chong et al.\ automated the generation of robot
simulation input for timber-frame assembly \cite{wong2022wood}, while Gao et
al.\ combined BIM-based task and motion planning for lightweight emergency
structures \cite{gao2022assembly}. These studies established that building
objects can be transformed into executable task information; they did not seek
to define a common acceptance test for the composed simulator stage.

Kim et al.\ developed a BIM-integrated planning and simulation system in which
building information was converted to an SDF world and used to configure robot
painting tasks \cite{kim2021bimrobot}. Zhu et al.\ subsequently proposed an
open, component-oriented pipeline that split IFC building elements, preserved
their identity, converted them into SDF smart components, and executed a
light-steel assembly sequence in ROS/Gazebo \cite{zhu2023ifcsdf}. This work
established automated IFC-to-simulator conversion as feasible, but assembly
interface validity and cross-component physical consistency remained outside
its scope.

Research then moved closer to action-level construction information. The RCAN
framework linked IFC-derived data to ordered construction actions and ROS
execution, enabling changes in component-oriented plans to be propagated to a
mobile manipulator \cite{zhu2024rcan}. Oyediran et al.\ integrated 4D BIM, a
construction-robot knowledge base, high-level task planning, and action-level
simulation for indoor wall-frame installation
\cite{oyediran2025fourdbim}. These systems enrich BIM with task knowledge and
make robot operations more context-aware. Their principal concern is the
creation and flow of planning information, rather than a formal test of
whether all assets, references, intermediate states, and physical assumptions
form a valid task scene.

Semantic world modelling provides a second route from BIM to robot execution.
Kim and Peavy represented building elements, spaces, and robot-relevant
relationships in a semantic building world \cite{kim2022semantic}. For mobile
robots, IFC has been converted into maps, waypoints, and schedule-aware Gazebo
environments \cite{moura2021localization,gopee2023navigation}. Knowledge-driven
inspection and progress-monitoring systems likewise use BIM to select targets
and interpret observations \cite{chen2023inspection,zhao2024progress}.
Ontologies have been proposed to regularize the vocabulary for navigation and
inspection \cite{bahreini2024ontology}, and semantic BIM attributes have been
used to alter path cost and clearance in safety-aware planning
\cite{amani2025pathfinding}. This line of work demonstrates the operational
value of semantics, but its main unit of analysis is a map, waypoint, or
knowledge base. Assembly readiness additionally requires valid contact
interfaces, physical parameters, object states, and task success conditions.

\subsection{Physics-ready assets and multi-scale scene composition}

Robotic simulation depends on asset properties that BIM models rarely provide
in executable form. Xu et al.\ addressed this issue for construction equipment
by fusing CAD geometry, equipment templates, and technical specifications into
articulated USD models \cite{xu2024physics}. Their method generated hierarchy,
joints, drives, materials, and collision representations for a wheel loader,
and used a large language model to extract unspecified mechanics. The study
substantially reduced modelling time and is a close precedent for automated
physics-ready generation. Its stated scope was equipment modelling rather than
environment modelling or equipment--environment interaction. It therefore
does not resolve the composition of building components, temporary states,
task interfaces, and robot capabilities into one assembly scene.

General-purpose robotics simulators make the representational gap visible in
different ways. MuJoCo emphasizes efficient rigid-body dynamics and
contact-rich model-based control \cite{todorov2012mujoco}; GPU-native
simulation environments such as Isaac Gym make it practical to execute many
robot worlds in parallel \cite{makoviychuk2021isaacgym}. These runtime
capabilities do not establish that a source model was compiled correctly.
Domain randomization can expose policies to variations in visual or physical
parameters \cite{tobin2017domain}, but it presupposes a valid nominal model and
defensible parameter ranges. It is therefore distinct from detecting a
reversed coordinate frame, a missing assembly mate, or a physical value with
no authoritative source.

OpenUSD is well suited to scene composition because its layers, references,
variants, and payloads separate asset definitions from scene opinions and
permit data to be loaded on demand \cite{openUSD}. This is important for
construction: a building shell may support navigation at coarse resolution
while an active connection requires detailed collision and contact geometry.
Ren et al.\ developed a specification-driven IFC-to-USD pipeline in which
information delivery requirements act as a scene contract for selecting and
mapping identifiers, hierarchy, geometry, materials, and attributes
\cite{ren2025ifcusd}. Their study establishes both lifecycle-oriented semantic
preservation and configurable IFC-to-USD delivery; neither can be claimed as
new here. What remains unresolved for robot construction is how a current task
should activate geometric scales, physical interfaces, temporary states, and
cross-asset acceptance conditions. IfcOpenShell provides access to IFC4.3
objects, property sets, spatial containment, and geometry
\cite{ifc43,ifcopenshellDocs}, but a domain-specific compilation policy is
still required to turn that information into an auditable task scene.

SimReady complements USD composition with machine-readable capabilities,
features, profiles, and validation rules
\cite{nvidiaSimReadyWorkflow}. Current profiles are valuable for checking
individual assets, including rigid-body, collider, material, hierarchy, and
robot-articulation requirements. An asset-level result cannot establish that a
particular robot, tool, work zone, component, and target interface have been
composed consistently. The official FAQ identifies scene-level specifications
and validators as a limitation of the current framework
\cite{nvidiaSimReadyFAQ}. \systemname{} therefore builds on, rather than
replaces, asset validation by adding relations that become meaningful only
after a construction task scene is composed.

Scene generation for embodied AI is advancing in parallel and narrows the
claim space further. SimRecon reconstructs compositional simulator scenes from
video through perception, object generation, and scene-graph-guided assembly
\cite{xia2026simrecon}. EmbodiedGen V2 generates editable, task-driven worlds
with affordances and cross-simulator deployment, and evaluates whether the
worlds support downstream policy learning \cite{wang2026embodiedgenv2}.
Agentic Real2Sim recovers geometry, physical parameters, robot interaction,
and object state from recorded episodes using vision-language agents
\cite{chen2026agenticreal2sim}. These systems show that agentic production of
runnable, physically enriched scenes is no longer an unoccupied research
problem. Their source of authority is visual observation or generative world
specification, whereas a construction scene must also remain accountable to
IFC identity, engineering interfaces, process state, and explicitly bounded
parameter sources. This difference motivates an auditable compiler and fault
benchmark rather than another general-purpose scene generator.

The use of a humanoid evaluation platform makes this cross-scale requirement
especially visible. Whole-body benchmarks show that locomotion and
manipulation outcomes depend on coupling among support, reach, contact, and
object dynamics \cite{sferrazza2024humanoidbench}. In a construction scene,
the coupling extends across building circulation, body placement, work-zone
occupancy, and millimetre-scale assembly geometry. A compiler must support
coarse global context and detailed active interfaces without representing
every building object at maximum fidelity. This is a scene-composition
requirement, not a claim that humanoids are universally preferable to wheeled
mobile manipulators. Nor is BIM-to-humanoid simulation itself new: Ye et al.\
convert BIM geometry and semantic information to a MuJoCo scene and demonstrate
task decomposition with Unitree G1 \cite{ye2026humanoiddt}. The present use of
G1 is therefore an evaluation stress test for task-scene readiness rather than
a standalone contribution.

\subsection{Closed-loop and robot-ready construction twins}

Digital-twin research has addressed the discrepancy between design models and
changing construction sites. Boje et al.\ argued for semantically structured
construction twins that connect heterogeneous representations and lifecycle
processes \cite{boje2020semantic}. Sacks et al.\ placed the twin within a
closed-loop production system in which observations and decisions continually
update construction control \cite{sacks2020digitaltwin}. These formulations
show that a design model is only one view of a changing physical system. They
do not, however, prescribe a simulator-specific readiness contract for a
robotic task.

Wang et al.\ proposed a closed-loop BIM-driven workflow in which a human could
revise task sequences, target positions, and motion plans when drywall
installation encountered as-designed/as-built differences
\cite{wang2024closeddt}. Tuomisto et al.\ used BIM to initialize LiDAR SLAM and
next-best-view planning for quadruped inspection, while allowing the map to
distinguish designed structures, missing elements, and dynamic obstacles
\cite{tuomisto2026inspection}. These approaches make robot operation robust to
observed site conditions, but they validate the world through runtime sensing
and interaction rather than through pre-runtime analysis of a composed
simulation scene.

The BIM2RDT work of Akhavian et al.\ further integrates BIM geometry, robot
observations, IoT data, semantic registration, and safety-aware navigation
\cite{akhavian2026bim2rdt}. It represents an important step toward dynamic,
robot-ready construction twins and LLM-orchestrated workflows. Its principal
technical contributions concern point-cloud registration, perception,
monitoring, and navigation constraints. The present work addresses a
complementary stage: compiling an auditable physical scene before simulation,
then determining whether its cross-asset relationships and construction-task
assumptions are internally ready for execution.

Engineering analysis can also serve as an acceptance gate. The integration of
BIM, finite-element analysis, and ROS has shown how temporary structural
stability can modify a robot assembly plan \cite{xu2025fem}. BIM-semantic-web
integration can supplement an observed component with design relationships and
requirements \cite{zhu2025perception}. These studies demonstrate that
readiness is conditional on the current state and action, not an intrinsic
label attached to an asset. \systemname{} adopts that conditional view while
limiting its present physical scope to rigid components and simplified contact
interfaces.

\subsection{Language models, agents, and engineering assurance}

Language models have begun to support construction-robot planning and
engineering-model generation. RoboGPT used ChatGPT to generate and revise
assembly sequences, demonstrating that natural-language reasoning can adapt a
robot plan across construction cases \cite{you2023robogpt}. Xu et al.\ used an
LLM more narrowly to extract equipment-joint parameters from technical
documents \cite{xu2024physics}; BIM2RDT uses language-model reasoning to assign
semantic priors and orchestrate parts of a robot-ready twin pipeline
\cite{akhavian2026bim2rdt}. These studies show the value of models that connect
heterogeneous descriptions and select procedural steps.

The wider robotics literature offers several patterns for grounding language
in executable behavior. SayCan combines language-model preferences with
learned affordance values \cite{ahn2022saycan}; Inner Monologue uses
environment feedback to revise embodied plans \cite{huang2022innermonologue}.
Code as Policies and ProgPrompt constrain generated output through executable
program structure or situated action specifications
\cite{liang2023code,singh2023progprompt}. Multimodal embodied models such as
PaLM-E and vision-language-action models such as RT-2 move perception and
action closer to the language model itself
\cite{driess2023palme,brohan2023rt2}. Multi-robot planning frameworks further
decompose tasks, allocate capabilities, and produce programmatic plans
\cite{kannan2023smartllm}. These methods strengthen task generation, but none
of them makes a language model an authoritative source for an engineering
mass, tolerance, coordinate frame, or IFC identity.

A plausible action sequence is not necessarily executable, and a plausible
numerical parameter is not necessarily authoritative. In \systemname{}, the
agent is not treated as a validator or a source of ground truth. It receives a
typed error report, chooses from bounded repair tools, records the evidence for
each mutation, and returns the scene to deterministic rules. This pattern
differs from direct language-to-scene generation because acceptance remains
independent of the model that proposed the repair. It also makes fixed-rule and
agentic methods comparable on identical injected faults.

The separation of proposal from acceptance is consistent with hybrid planning
in which a language model supplies a rough plan and a formal answer-set planner
establishes executability \cite{lin2025asp}. It also aligns with fault-oriented
robot assurance: injected faults can determine whether a digital
representation and its deployment process expose unsafe assumptions before
hardware execution \cite{frasheri2023safer}. In the present benchmark, the
object under test is the construction scene. The agent may select a repair,
but typed rules, source authority, and task-level checks determine whether that
repair is accepted.

\subsection{Synthesis of the research gap}

\Cref{tab:literature} compares the closest methodological precedents. The
reviewed literature establishes IFC-to-simulator conversion, BIM-derived robot
tasks, semantic maps, specification-driven IFC-to-USD delivery, physics-ready
USD asset generation, agentic simulation-world generation, BIM-to-humanoid
simulation, closed-loop twin updating, semantic perception, and
language-model-assisted planning. These capabilities are necessary parts of a
robot-ready workflow, but they operate on different units of analysis.
Conversion studies focus on delivered building information, physics-ready
studies on an equipment asset, generative-world studies on executable scenes,
planning studies on action information, and digital-twin studies on runtime
alignment with a site.

The literature search did not identify a study that jointly evaluates a
construction-task scene contract, task-activated multi-scale USD composition,
physical-property provenance, deterministic scene-level fault detection,
validator-grounded repair, and the downstream outcome of a humanoid assembly
task. This is deliberately narrower than a claim of first IFC--robot
interoperability or first USD validation; both are established areas. The
unresolved transition is from individually valid digital assets to an
auditable, composed, task-specific construction world.

\begin{table*}[t]
\centering
\caption{Capabilities of closely related BIM--robotics and physics-ready
modelling studies. A filled cell indicates that the capability was a
substantive part of the reported method or evaluation.}
\label{tab:literature}
\small
\resizebox{\textwidth}{!}{%
\begin{tabular}{lcccccccc}
\toprule
Study & BIM/IFC & Robot task & Physical asset & Multi-scale &
Scene rules & Provenance & Agent repair & Runtime outcome \\
\midrule
Kim et al.\ \cite{kim2021bimrobot}
& \checkmark & \checkmark & partial & -- & -- & -- & -- & \checkmark \\
Wong Chong et al.\ \cite{wong2022wood}
& \checkmark & \checkmark & partial & -- & -- & -- & -- & \checkmark \\
Zhu et al.\ \cite{zhu2023ifcsdf}
& \checkmark & \checkmark & partial & -- & -- & -- & -- & \checkmark \\
Oyediran et al.\ \cite{oyediran2025fourdbim}
& \checkmark & \checkmark & partial & -- & -- & -- & -- & \checkmark \\
Xu et al.\ \cite{xu2024physics}
& -- & -- & \checkmark & -- & asset & partial & partial & \checkmark \\
Ren et al.\ \cite{ren2025ifcusd}
& \checkmark & -- & partial & scene & specification & \checkmark & -- & -- \\
Wang et al.\ \cite{wang2024closeddt}
& \checkmark & \checkmark & partial & -- & runtime & -- & human & \checkmark \\
Gopee et al.\ \cite{gopee2023navigation}
& \checkmark & navigation & partial & map & runtime & -- & -- & \checkmark \\
Akhavian et al.\ \cite{akhavian2026bim2rdt}
& \checkmark & navigation & partial & -- & runtime & partial & partial & \checkmark \\
Ye et al.\ \cite{ye2026humanoiddt}
& \checkmark & \checkmark & partial & scene & -- & -- & -- & \checkmark \\
Wang et al.\ \cite{wang2026embodiedgenv2}
& -- & \checkmark & \checkmark & \checkmark & partial & partial & generation & \checkmark \\
\systemname{}
& \checkmark & \checkmark & \checkmark & \checkmark & scene &
\checkmark & \checkmark & \checkmark \\
\bottomrule
\end{tabular}
}
\end{table*}
